% Pursuing Stacks (cotes 134-2 to 134-8): archive cote 134-8, site page 5 (PDF page 6), typescript page 559.
% typescript. Draft compared: ps7.tex:195-261.
% AI-transcribed by Claude Opus 5.5 (claude-opus-5-5, adaptive thinking, effort high), Message Batches API, batch ps4 (msgbatch_012Aoc1NuskATPsyBeKSYK7J), 2026-11-12; not yet checked by eye.
\scanpage{134-8/5}
\mspage{559}

integrat or for \hand{$\Delta$}, the upper horizontal arrow in (7) is \uline{not} the
Kan-Dold-Puppe equivalence of categories, it has to be followed still
by the ``normalization'' operation. Thus \uline{we certainly should not expect} \margin{!x}
\uline{in any case the functor} (5) \uline{to be an equivalence} - \uline{however we suspect}
\uline{that when A is a test category} (maybe even a weak test category would
do it), \uline{then the lower horizontal arrow in} (7)
\[ (8) \qquad \mathrm{Hotab}_A \longrightarrow \mathrm{D.(Ab)} = \mathrm{H}_{\mathrm{o}}\mathrm{tab} \]
\uline{is an equivalence of categories}. \inserted[typed]{Whenever true, \xout{this}} \xout{\unsure{Suchxkindxofxstatement,}} for a given
category A , \inserted[typed]{this statement} looks like a reasonable substitute (on the level of the
relevant derived categories) \xout{of} for the Kan-Dold-Puppe theorem, known
in the two cases \hand{$\Delta$} and \hand{\unsure{$\bigcirc$}}.

There are however still two \inserted[typed]{important} extra features in the case A=\hand{$\Delta$}, which
deserve to be understood in the case of more general A . One is that
a map $u\colon F\to G$ in $A^{\wedge}_{ab}$ is in $W^{ab}_A$ (i.e. is an ``abelian weak equivalence''\hand{\unsure{)}}
iff it is in $W_A$ i.e. iff it is a weak equivalence, when forgetting the
abelian structures. In terms of a final object e in A (when such ob-%
ject exists indeed), viewing the categories $A_{/F}$ (for F in $A^{\wedge}_{ab}$) pointed \margin{it is enough that A be 1 connected instead of having a final object}
by the zero map $e \to F$ , which yields the necessary base-point \inserted{$e_F$} for
defining homotopy invariants $\pi_i(F)$ for $i \geqslant 0$, the relationship just
considered between abelian weak equivalence and weak equivalence, will
follow of course whenever we have functorial isomorphisms
\[ (9) \qquad \pi_i(F) \overset{\mathrm{def}}{=} \pi_i(A_{/F} \unsure{;|,} e_F) \simeq H_i(A,F) \xout{\illegible{1}} \quad , \]
which are known to exist indeed in the case A =\hand{$\Delta$}. It should be noted
that in section 92, when starting (in a somewhat casual way) with some
reflections on ``abelianization'', we introduced a category (denoted
by $\mathrm{Hotab}_A$) by localizing $A^{\wedge}_{ab}$ with respect to \inserted{the} maps which are weak
equivalences (when forgetting the abelian structures), whereas it has
now become clear that, in case \inserted[typed]{the latter} \xout{\unsure{thisxnotion}} should not coincide with
the \xout{\unsure{notion}} ``abelian weak equivalences'' defined in terms of linearization
or homology, it is the category (6) definitely which is the right one.
Still the question of defining isomorphisms (9) when A has a final
object, and whether the equality
\[ (10) \qquad W^{ab}_A = \xout{\illegible{1}} \ \mathrm{forg}_A{}^{-1}(W_A) \]
holds, where
\[ \mathrm{forg}_A\colon A^{\wedge}_{ab} \longrightarrow A^{\wedge} \]
is the ``forgetful functor'' , should be settled for general A.

\margin{5}
